\documentclass[12pt,oneside]{book}

%==================================================================
% METADATOS
%==================================================================
\newcommand{\universidad}{Universidad de Buenos Aires}
\newcommand{\facultad}{Facultad de Derecho}
\newcommand{\materia}{Derecho Constitucional}
\newcommand{\titulo}{El Recurso Extraordinario Federal}
\newcommand{\autor}{Isaac Edgar Camacho Ocampo}
\newcommand{\anio}{2026}

%==================================================================
% PAQUETES
%==================================================================
\usepackage[utf8]{inputenc}
\usepackage[T1]{fontenc}
\usepackage[spanish,es-nodecimaldot]{babel}

\usepackage[
    a4paper,
    top=2.5cm,
    bottom=2.5cm,
    left=3cm,
    right=2.5cm,
    headheight=15pt
]{geometry}

\usepackage{amsmath}
\usepackage{amssymb}
\usepackage{mathtools}

\usepackage{graphicx}
\usepackage{float}
\usepackage{caption}
\usepackage{subcaption}

\usepackage{booktabs}
\usepackage{array}
\usepackage{longtable}
\usepackage{tabularx}

\usepackage{enumitem}

\usepackage{xcolor}
\usepackage[most]{tcolorbox}

\usepackage{fancyhdr}
\usepackage{ragged2e}

\usepackage[
    colorlinks=true,
    linkcolor=blue!50!black,
    citecolor=green!40!black,
    urlcolor=blue!60!black,
    pdftitle={\titulo},
    pdfauthor={\autor},
    pdfsubject={Apuntes universitarios},
    bookmarks=true
]{hyperref}

%==================================================================
% CONFIGURACIÓN
%==================================================================
\graphicspath{
    {./img/}
    {./graficos/}
    {./figuras/}
}

\setcounter{tocdepth}{2}

\setlength{\parindent}{0pt}
\setlength{\parskip}{0.5em}

\setlist[itemize]{
    topsep=0.3em,
    itemsep=0.2em
}

\setlist[enumerate]{
    topsep=0.3em,
    itemsep=0.2em
}

\clubpenalty=10000
\widowpenalty=10000

%==================================================================
% COMANDOS
%==================================================================
\newcommand{\abs}[1]{\left|#1\right|}
\newcommand{\norm}[1]{\left\lVert#1\right\rVert}
\newcommand{\vect}[1]{\mathbf{#1}}
\newcommand{\deriv}[2]{\frac{d#1}{d#2}}
\newcommand{\pderiv}[2]{\frac{\partial#1}{\partial#2}}

%==================================================================
% ENTORNOS / CAJAS
%==================================================================
\newtcolorbox{definition}[1]{
    enhanced,
    breakable,
    title={Definición: #1},
    fonttitle=\bfseries,
    colback=blue!3,
    colframe=blue!50!black,
    boxrule=0.8pt,
    arc=2mm
}

\newtcolorbox{important}{
    enhanced,
    breakable,
    title={Importante},
    fonttitle=\bfseries,
    colback=yellow!8,
    colframe=orange!70!black,
    boxrule=0.8pt,
    arc=2mm
}

\newtcolorbox{example}[1]{
    enhanced,
    breakable,
    title={Ejemplo: #1},
    fonttitle=\bfseries,
    colback=green!4,
    colframe=green!40!black,
    boxrule=0.8pt,
    arc=2mm
}

\newtcolorbox{formula}[1]{
    enhanced,
    breakable,
    title={Fórmula / Regla: #1},
    fonttitle=\bfseries,
    colback=purple!4,
    colframe=purple!50!black,
    boxrule=0.8pt,
    arc=2mm
}

\newtcolorbox{exercise}[1]{
    enhanced,
    breakable,
    title={Ejercicio: #1},
    fonttitle=\bfseries,
    colback=gray!5,
    colframe=gray!60!black,
    boxrule=0.8pt,
    arc=2mm
}

\newtcolorbox{note}{
    enhanced,
    breakable,
    title={Nota},
    fonttitle=\bfseries,
    colback=white,
    colframe=gray!50,
    boxrule=0.6pt,
    arc=2mm
}

%==================================================================
% CONFIGURACIÓN FINAL
%==================================================================
\pagestyle{fancy}
\fancyhf{}

\fancyhead[L]{\nouppercase{\leftmark}}
\fancyhead[R]{\materia}

\fancyfoot[C]{\thepage}

\renewcommand{\headrulewidth}{0.4pt}
\renewcommand{\footrulewidth}{0pt}

\fancypagestyle{plain}{
    \fancyhf{}
    \fancyfoot[C]{\thepage}
    \renewcommand{\headrulewidth}{0pt}
}

%==================================================================
% DOCUMENTO
%==================================================================
\begin{document}

\frontmatter

%------------------------------------------------------------------
% PORTADA
%------------------------------------------------------------------
\begin{titlepage}

\thispagestyle{empty}

\begin{center}

\vspace*{1cm}

{\large \textsc{\universidad}}

\vspace{0.2cm}

{\large \textsc{\facultad}}

\vspace{3cm}

{\Huge \textbf{\materia}}

\vspace{1cm}

{\LARGE \titulo}

\vspace{0.8cm}

\textit{El estudio constante convierte la comprensión en conocimiento.}

\vspace{1.2cm}

\rule{0.70\textwidth}{0.5pt}

\vspace{0.5cm}

{\large Apuntes de estudio}

\vspace{0.5cm}

\rule{0.70\textwidth}{0.5pt}

\vfill

{\large \autor}

\vspace{0.3cm}

{\large \anio}

\end{center}

\vfill

\begin{flushleft}
\textit{Este es un modesto aporte a los estudiantes de ``\materia'' no pretende sustituir el material oficial de la materia.}
\end{flushleft}

\end{titlepage}

%------------------------------------------------------------------
% ÍNDICE
%------------------------------------------------------------------
\tableofcontents

\mainmatter

%==================================================================
% CAPÍTULO 1
%==================================================================
\chapter[Concepto y requisitos formales]{Concepto del recurso extraordinario federal y requisitos formales}

\section{Visión general del tema}

\subsection{Temas tratados}

El estudio del recurso extraordinario federal (REF) comprende los siguientes temas:

\begin{itemize}
    \item El concepto y la estructura del REF como recurso de impugnación, y no como instancia de jurisdicción originaria.
    \item Los requisitos formales: la Acordada 4/2007 de la Corte Suprema (extensión, carátula), el plazo de 10 días y la reserva del caso federal.
    \item Los requisitos comunes a todo recurso judicial: la existencia de un gravamen actual y de un proceso judicial en trámite.
    \item Los requisitos propios del REF: sentencia definitiva o asimilable, cuestión federal (simple, compleja directa e indirecta), decisión contraria al derecho federal invocado y relación directa e inmediata.
    \item La queja por recurso extraordinario denegado y el instituto del \textit{per saltum}.
    \item La jurisprudencia reciente de control de constitucionalidad: DNU 70/2023, reforma laboral, ley de medios y financiamiento universitario.
\end{itemize}

\subsection{Hilo conductor}

Todo el tema puede pensarse como el manual de instrucciones para acudir, una sola vez y de forma correcta, a la Corte Suprema. En primer lugar, hay que cumplir el protocolo de ingreso: los \textbf{requisitos formales}, comparables con completar bien un formulario antes de entrar a un edificio público. En segundo lugar, hay que demostrar que existe un motivo real para acudir: los \textbf{requisitos comunes}, comparables con tener realmente un problema que resolver. Por último, hay que probar que el problema es precisamente del tipo que la Corte atiende, es decir, uno que involucra la Constitución: los \textbf{requisitos propios}, esto es, la cuestión federal.

Si en algún paso de este recorrido algo falla, la puerta no se abre y solo queda un último recurso: la queja.

\subsection{Relevancia profesional y formativa}

El recurso extraordinario federal es la herramienta procesal que permite llevar un caso hasta la Corte Suprema cuando está en juego la supremacía de la Constitución. Dominar sus requisitos resulta indispensable para cualquier abogado litigante: un solo error de forma (un plazo vencido, una página de más) puede cerrar para siempre el acceso al máximo tribunal, sin importar cuán sólido sea el argumento de fondo.

Comprender la diferencia entre cuestión federal simple y compleja permite anticipar el destino de cualquier norma cuestionada (un DNU, una reforma laboral, una ley provincial) y asesorar con criterio a un cliente sobre sus verdaderas posibilidades de éxito.

Su utilidad excede además el ámbito profesional: enseña a distinguir lo esencial de lo accesorio, a verificar los requisitos antes de invertir esfuerzo en el fondo de una cuestión y a comprender cómo funciona el control de constitucionalidad que protege los derechos de todos los ciudadanos frente al Estado.

\subsection{Mapa de contenidos}

La siguiente tabla ordena los temas según los cuatro bloques en que se organiza el estudio, junto con la pregunta que orienta cada uno.

\begin{longtable}{@{}>{\raggedright\arraybackslash}p{0.60\textwidth}>{\raggedright\arraybackslash}p{0.35\textwidth}@{}}
\toprule
\textbf{Tema} & \textbf{Pregunta orientadora} \\
\midrule
\endfirsthead
\toprule
\textbf{Tema} & \textbf{Pregunta orientadora} \\
\midrule
\endhead
\bottomrule
\endfoot
\bottomrule
\endlastfoot

\multicolumn{2}{@{}l}{\textbf{Bloque I. Requisitos formales del recurso}} \\
\addlinespace[0.2em]
El recurso extraordinario como mecanismo de impugnación: concepto, recurso frente a jurisdicción originaria. &
¿Qué es y cuándo procede el REF? \\
\addlinespace[0.3em]
Acordada 4/2007 de la Corte Suprema: extensión, carátula, formato del escrito. &
¿Qué límites formales fija la Corte? \\
\addlinespace[0.3em]
Plazo y tribunal de interposición: 10 días, mismo tribunal, tempestividad. &
¿Cuándo y ante quién se presenta? \\
\addlinespace[0.3em]
El planteo del caso federal en cada instancia: reserva, mantenimiento, oportunidad. &
¿Cómo se preserva la vía federal? \\
\midrule

\multicolumn{2}{@{}l}{\textbf{Bloque II. Requisitos comunes a todo recurso}} \\
\addlinespace[0.2em]
Gravamen actual y existencia de un proceso judicial: perjuicio concreto, tribunal judicial. &
¿Qué exige cualquier recurso, sin excepción? \\
\midrule

\multicolumn{2}{@{}l}{\textbf{Bloque III. Requisitos propios del recurso extraordinario}} \\
\addlinespace[0.2em]
Sentencia definitiva o asimilable: medidas cautelares, perjuicio irreparable. &
¿Qué resoluciones habilitan el recurso? \\
\addlinespace[0.3em]
Cuestión federal simple: interpretación de norma federal. &
¿Qué alcance tiene la norma federal? \\
\addlinespace[0.3em]
Cuestión federal compleja, directa e indirecta: conflicto de normas, jerarquía constitucional. &
¿Qué normas están en pugna y en qué nivel? \\
\addlinespace[0.3em]
Decisión contraria al derecho federal y relación directa e inmediata: conexión entre la cuestión federal y el fallo. &
¿Por qué la sentencia vulnera el derecho federal invocado? \\
\midrule

\multicolumn{2}{@{}l}{\textbf{Bloque IV. La queja y casos de aplicación reciente}} \\
\addlinespace[0.2em]
La queja por recurso extraordinario denegado: depósito, vía directa ante la Corte. &
¿Qué hacer si se rechaza el REF? \\
\addlinespace[0.3em]
Jurisprudencia y control de constitucionalidad reciente: DNU 70/2023, reforma laboral, \textit{per saltum}, ley de medios, financiamiento universitario. &
¿Cómo se aplican estos conceptos en casos reales? \\
\end{longtable}

\section{El recurso extraordinario como mecanismo de impugnación}

El recurso extraordinario es, ante todo, un \textbf{recurso}, es decir, un mecanismo de impugnación que, por su propia naturaleza, presupone que ya intervino un tribunal con anterioridad.

\begin{definition}{Recurso extraordinario}
Recurso que se interpone en el marco de un juicio que ya está en curso y habiendo intervenido otro tribunal.
\end{definition}

Por lo tanto, no se está ante la denominada \textbf{jurisdicción originaria} de la Corte Suprema, sino ante una instancia recursiva posterior a la intervención de otro tribunal.

\section{Clasificación de los requisitos}

Para precisar la consistencia y la definición del recurso extraordinario, el estudio se organiza en tres tipos de requisitos: formales, comunes y propios. Esta división tiene una finalidad metodológica: ayudar a comprender qué es, en definitiva, el recurso extraordinario.
% TODO: contenido incompleto o ambiguo en las notas originales (la frase sobre la preferencia por "nombres y divisiones" aparece truncada o con posible error de transcripción)

\begin{definition}{Tipos de requisitos del recurso extraordinario}
\begin{itemize}
    \item \textbf{Requisitos formales:} cuestiones objetivas o de forma, comunes a la presentación del recurso.
    \item \textbf{Requisitos comunes:} los presentes en cualquier tipo de recurso judicial, no solo en el extraordinario.
    \item \textbf{Requisitos propios:} lo particular del recurso extraordinario; lo esencial para saber cuándo resulta procedente interponerlo.
\end{itemize}
\end{definition}

\begin{important}
Cualquier inobservancia de los requisitos formales conduce, en la generalidad de los casos, al rechazo del recurso.

Existe una excepción: el supuesto de quien litiga \textit{in forma pauperis}, es decir, cuando resulta muy complicado exigirle a un ciudadano el cumplimiento estricto de estos requisitos por encontrarse en una situación particular de vulnerabilidad o de imposibilidad material de cumplirlos.
\end{important}

\section{Requisitos formales}

\subsection{La Acordada 4/2007 de la Corte Suprema}

Algunos requisitos formales (la extensión de las páginas, el interlineado y la forma de redacción del escrito) no están regulados en el Código Procesal Civil y Comercial, sino en una acordada de la Corte Suprema: la \textbf{Acordada 4 del año 2007}. Esta acordada de la Corte Suprema de Justicia de la Nación reglamenta las condiciones formales de interposición del recurso extraordinario federal y de la queja por recurso extraordinario denegado.

En su momento la acordada fue cuestionada, porque parte de la doctrina entendía que la Corte carecía de una función o competencia regulatoria de orden procesal. Sin embargo, existe cierta costumbre, sobre todo en el derecho americano, conforme a la cual la Corte regula algunas condiciones formales y procesales del trámite ante sí misma. Esa recepción del derecho americano, sumada a la idea de que la Corte tiene capacidad para regular algunos aspectos del proceso más allá de lo establecido en el Código Procesal, ha llevado a que la Acordada 4/2007 esté totalmente aceptada y sea muy difícil de impugnar.

Mediante esta acordada la Corte buscó estratificar y fijar límites a la extensión de los recursos, para evitar que resulten engorrosos, organizar su propio trabajo y ganar eficiencia en el análisis de este tipo de presentaciones. Estos criterios procuran racionalizar y sistematizar la forma en que se presentan los escritos ante la Corte. La acordada constituye una facultad de la Corte para dictar actos generales, como el reglamento de las condiciones de interposición del recurso extraordinario (designado coloquialmente «Rex»).

\subsubsection*{Requisitos establecidos por la Acordada 4/2007}

\begin{table}[H]
\centering
\begin{tabularx}{\textwidth}{@{}>{\raggedright\arraybackslash}p{0.38\textwidth}>{\raggedright\arraybackslash}X@{}}
\toprule
\textbf{Aspecto} & \textbf{Requisito} \\
\midrule
Extensión máxima del recurso extraordinario & 40 páginas. \\
Extensión máxima de la queja & 10 páginas. \\
Renglones por página & Máximo de 26 renglones por página. \\
Tamaño de letra & Mínimo: cuerpo 12. \\
Carátula & Debe completarse con los datos del expediente y un resumen del proceso. \\
\bottomrule
\end{tabularx}
\end{table}

\subsection{Plazo de interposición y tribunal ante el cual se presenta}

\begin{itemize}
    \item \textbf{Plazo:} 10 días hábiles desde la notificación de la resolución que se impugna.
    \item \textbf{Tribunal:} el recurso debe interponerse ante el mismo tribunal que dictó la resolución impugnada.
    \item \textbf{Condición del tribunal:} debe ser el tribunal superior de la causa, es decir, el máximo tribunal al que se pueda llegar antes de la Corte, sin que existan otros medios recursivos disponibles.
\end{itemize}

\begin{example}{Tribunal de interposición en el ámbito provincial}
Si el tribunal que dictó la resolución impugnada es la Corte Suprema de una provincia y corresponde recurrir a la Corte Federal, el recurso extraordinario se interpone ante ese mismo Tribunal Superior provincial. Como ejemplo se mencionó el Tribunal Superior de Chaco, en tanto máxima autoridad judicial dentro de esa provincia.
\end{example}

\textbf{Oportunidad del planteo.} La cuestión federal debe haber sido introducida oportunamente: es necesario que ya haya sido planteada con anterioridad en el proceso, y no puede surgir recién en la última instancia, cuando la parte advierte que está perdiendo el litigio. De lo contrario, el planteo sería el fruto de una reflexión tardía del procedimiento o del proceso judicial.

\subsection{El tribunal superior de la causa en el ámbito provincial y en la CABA}

En materia de derecho común, es decir, derecho ordinario, el último tribunal previo a la Corte Suprema federal es, en el ámbito de una provincia, la \textbf{Corte Suprema de la Provincia}. Debe agotarse esa instancia y recién después, eventualmente, plantearse el recurso extraordinario, siempre que exista una cuestión federal.

Este mismo razonamiento se aplica en la Ciudad Autónoma de Buenos Aires (CABA), conforme al precedente identificado como «Levinas». En materia de derecho ordinario o común (el de las cámaras nacionales), antes de acudir a la Corte Suprema debe pasarse por el \textbf{Tribunal Superior de Justicia de la Ciudad Autónoma}, por ser el máximo tribunal de la jurisdicción en la que se litiga. De ahí que deba recurrirse primero ante ese Tribunal Superior de Justicia de la CABA.

\subsection{El planteo del caso federal en cada instancia («reserva»)}

Para que proceda el recurso extraordinario, la parte debe haber efectuado la reserva del caso federal desde el primer escrito y en cada instancia relevante del proceso.

En la jerga, la reserva del caso federal consiste, de forma muy sucinta, en manifestar la reserva de apelar ante la Corte por encontrarse afectado, por ejemplo, el derecho de propiedad o las garantías del debido proceso.

\begin{definition}{Planteo del caso federal (reserva)}
Advertir al tribunal que se está en presencia de un tema que suscita una cuestión federal, determinarla y dar la fundamentación de los artículos de la Constitución y de los precedentes involucrados, enumerándolos aunque sea de forma sucinta.
\end{definition}

\begin{note}
Se señaló que la expresión «reserva» es técnicamente cuestionable, porque los derechos no se reservan. La observación es técnicamente correcta; no obstante, suelen emplearse los términos aceptados en la jerga y en la actividad profesional para no generar confusión. En rigor, no se trata de guardar algo para más adelante: es más que una reserva, es el \textbf{planteo efectivo del caso federal}.
\end{note}

Este planteo debe reiterarse en cada instancia nueva del proceso: en primera instancia, al interponer o contestar la demanda; en segunda instancia, al expresar agravios.

\begin{important}
En cada instancia judicial debe hacerse un planteo efectivo de qué se entiende por caso federal y por qué se entiende que la causa conlleva un caso federal.
\end{important}

\clearpage
\section{Cuadro Cornell}

\begin{longtable}{@{}>{\raggedright\arraybackslash}p{0.27\textwidth}>{\raggedright\arraybackslash}p{0.68\textwidth}@{}}
\toprule
\textbf{Preguntas / palabras clave} & \textbf{Notas / respuestas} \\
\midrule
\endfirsthead
\toprule
\textbf{Preguntas / palabras clave} & \textbf{Notas / respuestas} \\
\midrule
\endhead
\endfoot
\endlastfoot

\textbf{¿Qué es el recurso extraordinario federal y por qué no es jurisdicción originaria?} &
Es un mecanismo de impugnación que se interpone en el marco de un juicio ya en curso y habiendo intervenido otro tribunal. Por eso presupone la intervención previa de otro tribunal y no constituye jurisdicción originaria de la Corte Suprema. \\
\addlinespace[0.6em]

\textbf{¿En qué tipos de requisitos se organiza su estudio?} &
En requisitos formales (cuestiones objetivas o de forma), requisitos comunes (presentes en cualquier recurso judicial) y requisitos propios (lo particular del recurso extraordinario, esencial para saber cuándo procede). \\
\addlinespace[0.6em]

\textbf{¿Qué consecuencia tiene inobservar los requisitos formales? ¿Hay alguna excepción?} &
En la generalidad de los casos, el rechazo del recurso. La excepción es quien litiga \textit{in forma pauperis}, cuando exigirle el cumplimiento estricto resulta muy complicado por su situación de vulnerabilidad o imposibilidad material. \\
\addlinespace[0.6em]

\textbf{¿Qué regula la Acordada 4/2007 y por qué se la considera aceptada?} &
Reglamenta las condiciones formales de interposición del REF y de la queja. Fue cuestionada porque parte de la doctrina negaba a la Corte competencia regulatoria procesal, pero hoy está totalmente aceptada y es muy difícil de impugnar, por la recepción del derecho americano y la idea de que la Corte puede regular aspectos del proceso. \\
\addlinespace[0.6em]

\textbf{¿Qué límites formales fija la Acordada 4/2007?} &
Extensión máxima de 40 páginas para el recurso y de 10 para la queja; máximo de 26 renglones por página; letra mínima de cuerpo 12; y obligación de completar una carátula con los datos del expediente y un resumen del proceso. \\
\addlinespace[0.6em]

\textbf{¿Cuál es el plazo para interponer el REF y ante qué tribunal?} &
10 días hábiles desde la notificación de la resolución impugnada, ante el mismo tribunal que la dictó, que debe ser el tribunal superior de la causa, sin otros medios recursivos disponibles. \\
\addlinespace[0.6em]

\textbf{¿Cuál es el tribunal superior de la causa en derecho común?} &
En una provincia, la Corte Suprema provincial. En la CABA, conforme al precedente «Levinas», el Tribunal Superior de Justicia de la Ciudad Autónoma. \\
\addlinespace[0.6em]

\textbf{¿Qué implica que el planteo sea oportuno?} &
Que la cuestión federal ya haya sido planteada con anterioridad en el proceso; no puede surgir en la última instancia como fruto de una reflexión tardía. \\
\addlinespace[0.6em]

\textbf{¿Qué es la «reserva» del caso federal y cuándo debe hacerse?} &
Técnicamente es el planteo efectivo de la cuestión federal: advertir la cuestión, determinarla y fundamentarla con los artículos constitucionales y precedentes involucrados, aunque sea de forma sucinta. Debe hacerse desde el primer escrito y reiterarse en cada instancia (en primera, al interponer o contestar la demanda; en segunda, al expresar agravios). \\

\midrule
\multicolumn{2}{@{}p{\dimexpr0.95\textwidth+2\tabcolsep\relax}@{}}{%
\textbf{Resumen:} El REF es un mecanismo de impugnación que presupone la intervención previa de otro tribunal. Su estudio se divide en requisitos formales, comunes y propios. Los formales (Acordada 4/2007, plazo de 10 días hábiles, tribunal de interposición, oportunidad y planteo del caso federal en cada instancia) deben cumplirse estrictamente, pues su inobservancia conduce, en general, al rechazo del recurso.
} \\
\bottomrule
\end{longtable}

%==================================================================
% CAPÍTULO 2
%==================================================================
\chapter[Requisitos comunes a todo recurso]{Requisitos comunes a todo recurso}

Los requisitos comunes son aquellos que hacen a la obligatoriedad presente en cualquier tipo de recurso judicial, y no solo en los recursos extraordinarios.

\section{El gravamen o perjuicio actual}

La existencia de un perjuicio es un presupuesto clave para poder recurrir.

\begin{definition}{Gravamen}
Perjuicio acreditado en cabeza de quien impugna. La decisión cuestionada debe generar, inevitablemente, un perjuicio concreto a quien recurre.
\end{definition}

\begin{important}
El daño debe ser \textbf{actual}: no puede ser conjetural ni futuro sin demasiado anclaje temporal. Debe ser presente o inmediato.
\end{important}

\section{La existencia de un proceso judicial}

Debe existir un juicio, es decir, un proceso judicial, en contraposición, por ejemplo, a una instancia administrativa. Aunque pueda parecer obvio, este punto es un presupuesto clave: la existencia de un proceso judicial es requisito ineludible para la procedencia de cualquier recurso, incluido el extraordinario.

\clearpage
\section{Cuadro Cornell}

\begin{longtable}{@{}>{\raggedright\arraybackslash}p{0.27\textwidth}>{\raggedright\arraybackslash}p{0.68\textwidth}@{}}
\toprule
\textbf{Preguntas / palabras clave} & \textbf{Notas / respuestas} \\
\midrule
\endfirsthead
\toprule
\textbf{Preguntas / palabras clave} & \textbf{Notas / respuestas} \\
\midrule
\endhead
\endfoot
\endlastfoot

\textbf{¿Qué son los requisitos comunes?} &
Los presentes en cualquier tipo de recurso judicial, no solo en el extraordinario. \\
\addlinespace[0.6em]

\textbf{¿Qué es el gravamen y cómo debe ser?} &
Es el perjuicio acreditado en cabeza de quien impugna: la decisión cuestionada debe generarle un perjuicio concreto. Debe ser actual, no conjetural ni futuro sin demasiado anclaje temporal; es decir, presente o inmediato. \\
\addlinespace[0.6em]

\textbf{¿Por qué se exige un proceso judicial?} &
Porque la existencia de un juicio o proceso judicial (y no, por ejemplo, una instancia administrativa) es un presupuesto clave e ineludible para la procedencia de cualquier recurso, incluido el extraordinario. \\

\midrule
\multicolumn{2}{@{}p{\dimexpr0.95\textwidth+2\tabcolsep\relax}@{}}{%
\textbf{Resumen:} Los requisitos comunes a todo recurso judicial son la existencia de un gravamen o perjuicio actual, concreto y acreditado en cabeza de quien recurre, y la existencia de un proceso judicial en trámite.
} \\
\bottomrule
\end{longtable}

%==================================================================
% CAPÍTULO 3
%==================================================================
\chapter[Requisitos propios del recurso extraordinario]{Requisitos propios del recurso extraordinario}

Los requisitos propios constituyen el núcleo más sustancial del recurso extraordinario: aquello que permite distinguirlo de cualquier otro recurso judicial, es decir, lo que hace al recurso extraordinario y que otros recursos no tienen.

\section{Sentencia definitiva o asimilable}

\begin{definition}{Sentencia definitiva o asimilable}
El recurso extraordinario, en principio, no está previsto para actos interlocutorios ni para medidas cautelares. Debe tratarse de sentencias definitivas o asimilables a definitivas, es decir, aquellas que ponen fin al proceso.

Una resolución es \textbf{asimilable a definitiva} cuando ya no hay marcha atrás, o en todo caso solo hay marcha adelante; es decir, cuando puede generarse un perjuicio irreparable.
\end{definition}

\subsection{Las medidas cautelares y su carácter provisional}

Las medidas cautelares son medidas provisionales, dictadas antes o en el curso del proceso, cuyo objeto fundamental es \textbf{asegurar el objeto del proceso}: que aquello que se resuelva en última instancia pueda efectivamente cumplirse, evitando que el mero paso del tiempo genere una situación irreversible.

Se planteó que algunos consideran a las medidas cautelares como sentencias anticipadas. % TODO: contenido incompleto o ambiguo en las notas originales (la intervención del alumno "No se recupera de cobrar" aparece incompleta)
Sin embargo, aunque siempre hay matices grises y resulta muy difícil diferenciar lo conceptual de lo que ocurre en la realidad, \textbf{técnicamente no son sentencias}, porque cualquier medida cautelar debe ser luego ratificada mediante la decisión de fondo.

\begin{example}{Medida cautelar que ordena proveer una vacuna a un menor}
Si se obtiene una cautelar que ordena proveer una vacuna a un menor, no puede darse por terminado el juicio con el argumento de que ya se obtuvo lo pretendido. Esa actitud puede resultar nociva: aunque nadie vaya a retrotraer la situación, es posible que luego se reclamen daños y perjuicios.
\end{example}

\subsection{Excepciones: cuando la Corte revisa medidas cautelares}

Excepcionalmente, en casos de especial trascendencia, de gravedad institucional o que involucran pluralidad de sujetos con un fuerte impacto social, la Corte se aboca a recursos extraordinarios interpuestos contra medidas cautelares, en la medida en que puedan generar perjuicios irreparables.
% TODO: contenido incompleto o ambiguo en las notas originales (el texto original dice "casos de norma trascendencia")

\begin{important}
No cualquier medida cautelar será abierta por la Corte: deben ser casos que procuren un daño irreparable o en los que haya muchas partes interesadas.
\end{important}

\begin{example}{Grupo Clarín y la Ley de Medios}
Cuando se sancionó la ley de medios, el grupo empresario la impugnó y, además, interpuso una medida cautelar para suspender la aplicación de la sanción hasta que se resolviera la validez de la nueva ley. La Corte abordó el recurso extraordinario pese a estar dirigido contra una medida cautelar, debido a que se trataba de una ley recientemente dictada que había generado un enorme planteo y un fuerte efecto social.
\end{example}

\begin{example}{Ley de financiamiento de la universidad pública}
En este caso no se planteó la invalidez de la norma, sino que la universidad sostuvo que no se estaba cumpliendo con la ley de financiamiento; es decir, se discutía si se cumplía o no con la obligación legal de proveer recursos a la universidad, a todas sus dependencias y a su personal docente y no docente. El fuero contencioso administrativo hizo lugar a la medida cautelar, pero le dio al Estado la posibilidad de que la cuestión fuera revisada en última instancia por la Corte, dada la delicadeza política y las múltiples implicancias del caso.

A los fines del recurso extraordinario, lo que se discutía en ese caso era la habilitación de la vía extraordinaria contra una medida cautelar.
\end{example}

\section{Cuestión federal}

En todo recurso extraordinario debe demostrarse la existencia de una controversia jurídica susceptible de ser resuelta por los tribunales, en la que esté en juego un caso o cuestión federal.

Según lo consignado en las notas, la Ley 48 establece la distinción entre cuestión federal simple y cuestión federal compleja a los fines de la procedencia del recurso extraordinario.

\subsection{Cuestión federal simple}

\begin{definition}{Cuestión federal simple}
Controversia en la que está en tela de juicio la interpretación de la ley, de la Constitución o de un tratado. Lo que se discute es el real alcance de la norma federal (entendida como la Constitución, las normas federales dictadas por el Congreso o los tratados internacionales), sin que exista un conflicto entre dos normas.
\end{definition}

\begin{example}{Alcance de la «indemnización plena»}
Si para resolver el caso está en juego desentrañar qué alcance corresponde darle al concepto de «indemnización plena», es decir, el sentido de esa expresión normativa, se trata de una cuestión federal simple.
\end{example}

\subsection{Cuestión federal compleja: directa e indirecta}

\begin{definition}{Cuestión federal compleja}
A diferencia de la cuestión federal simple, en la que se busca dar el verdadero alcance a los términos de una norma, en la cuestión federal compleja existe un \textbf{conflicto normativo} entre dos o más normas.
\end{definition}

\begin{table}[H]
\centering
\begin{tabularx}{\textwidth}{@{}>{\raggedright\arraybackslash}p{0.22\textwidth}>{\raggedright\arraybackslash}X>{\raggedright\arraybackslash}X@{}}
\toprule
\textbf{Tipo} & \textbf{Normas en conflicto} & \textbf{Ejemplo} \\
\midrule
Compleja directa &
El conflicto se da entre la Constitución Nacional y una norma inferior. &
Una ley o un decreto que contradicen directamente una disposición constitucional. \\
\addlinespace[0.4em]
Compleja indirecta &
El conflicto se da entre dos o más normas, todas ellas inferiores a la Constitución. &
Una norma provincial que controvierte una norma del Congreso; una norma provincial que controvierte el alcance de un tratado internacional. \\
\bottomrule
\end{tabularx}
\caption{Cuestión federal compleja directa e indirecta}
\end{table}

En la cuestión federal compleja indirecta la Constitución ya no está en tela de juicio: lo que se discute son controversias entre dos o más normas inferiores.

\subsubsection*{Los tratados del artículo 75, inciso 22}

El artículo 75, inciso 22 de la Constitución Nacional otorga jerarquía constitucional a los tratados internacionales de derechos humanos allí enumerados y prevé el mecanismo por el cual otros tratados pueden alcanzar dicha jerarquía.

Los tratados del artículo 75, inciso 22 son parte de la Constitución y, a los efectos de una cuestión federal, se los trata como tal. Por lo tanto:

\begin{itemize}
    \item Si se trata de un tratado incorporado a la Constitución por el mecanismo que prevé el inciso 22, la controversia constituye una \textbf{cuestión federal compleja directa}.
    \item Si la controversia se da entre un tratado internacional que no tiene estatus constitucional (por ejemplo, el tratado del Mercosur) y una ley nacional, existe una controversia entre dos normas inferiores a la Constitución y, por lo tanto, se trata de una \textbf{cuestión federal compleja indirecta}.
\end{itemize}

La finalidad de este requisito es que el recurso extraordinario asegure, en estos casos, el orden de prevalencia y de preeminencia de las normas. La Constitución es la norma superior en jerarquía, pero también establece cómo se ordenan las jerarquías posteriores: un tratado es superior a una ley. Si hubiera una discusión en ese plano, puede existir una cuestión federal compleja indirecta, que permite asegurar el orden de prelación que las normas tienen después de la Constitución.

\section{Decisión contraria al derecho federal invocado}

\begin{definition}{Decisión contraria al derecho federal invocado}
Para que el recurso extraordinario sea procedente, la decisión impugnada debe ser contraria al derecho federal invocado por quien recurre. Para ser pertinente y admitirse, el recurso debe demostrar esa contradicción.
\end{definition}

\begin{example}{Norma provincial y norma federal en pugna}
Si en el proceso estaban en pugna una norma provincial y una norma federal, y la instancia anterior resolvió que debía prevalecer la norma federal, no hay allí una cuestión constitucional en juego, porque la propia Constitución prevé que las normas federales son superiores a las locales. En ese supuesto no existe contradicción con el derecho federal invocado.
\end{example}

\begin{example}{Financiamiento educativo}
En el caso del financiamiento educativo estaban en juego, a priori, la autonomía universitaria y el derecho de asistencia de profesores y alumnos, entre otros temas que pueden destacarse en el marco de una cuestión federal.
\end{example}

\section{Relación directa e inmediata}

\begin{definition}{Relación directa e inmediata}
La resolución de la cuestión federal debe tener relación directa con la decisión del juicio. No puede tratarse de una controversia meramente conjetural o académica: la controversia judicial debe estar íntimamente relacionada con la cuestión federal invocada.
\end{definition}

\section{Síntesis de los requisitos propios}

\begin{table}[H]
\centering
\begin{tabularx}{\textwidth}{@{}>{\raggedright\arraybackslash}p{0.25\textwidth}>{\raggedright\arraybackslash}X@{}}
\toprule
\textbf{Cuestión federal} & \textbf{En qué consiste} \\
\midrule
Simple & Interpretación de una ley federal, de la Constitución o de un tratado (alcance de la norma). \\
\addlinespace[0.3em]
Compleja directa & Conflicto entre la Constitución Nacional y una norma inferior. \\
\addlinespace[0.3em]
Compleja indirecta & Conflicto entre dos o más normas, todas inferiores a la Constitución. \\
\bottomrule
\end{tabularx}
\caption{Cuestión federal simple y compleja}
\end{table}

\clearpage
\section{Cuadro Cornell}

\begin{longtable}{@{}>{\raggedright\arraybackslash}p{0.27\textwidth}>{\raggedright\arraybackslash}p{0.68\textwidth}@{}}
\toprule
\textbf{Preguntas / palabras clave} & \textbf{Notas / respuestas} \\
\midrule
\endfirsthead
\toprule
\textbf{Preguntas / palabras clave} & \textbf{Notas / respuestas} \\
\midrule
\endhead
\endfoot
\endlastfoot

\textbf{¿Qué son los requisitos propios del REF?} &
Son lo particular del recurso extraordinario, aquello que lo distingue de cualquier otro recurso judicial y permite saber cuándo es procedente interponerlo. \\
\addlinespace[0.6em]

\textbf{¿Qué tipo de resoluciones habilita el REF?} &
Sentencias definitivas o asimilables a definitivas: las que ponen fin al proceso o generan un perjuicio irreparable si no se revisan (cuando ya no hay marcha atrás). En principio no se prevé para actos interlocutorios ni medidas cautelares. \\
\addlinespace[0.6em]

\textbf{¿Qué son las medidas cautelares y por qué no son sentencias?} &
Son medidas provisionales cuyo objeto es asegurar que lo resuelto en última instancia pueda cumplirse. Técnicamente no son sentencias, porque deben ser ratificadas luego mediante la decisión de fondo. \\
\addlinespace[0.6em]

\textbf{¿Pueden recurrirse medidas cautelares mediante REF?} &
Excepcionalmente, cuando hay gravedad institucional, pluralidad de sujetos o fuerte impacto social, y en la medida en que puedan generar perjuicios irreparables. Casos citados: Grupo Clarín y Ley de Medios; financiamiento universitario. \\
\addlinespace[0.6em]

\textbf{¿Qué es una cuestión federal simple?} &
La controversia sobre el real alcance o interpretación de una norma federal (Constitución, ley del Congreso o tratado), sin conflicto entre normas. \\
\addlinespace[0.6em]

\textbf{¿Qué es una cuestión federal compleja directa?} &
El conflicto normativo se da directamente entre la Constitución Nacional y una norma inferior. Los tratados con jerarquía constitucional del artículo 75, inciso 22 se tratan como parte de la Constitución. \\
\addlinespace[0.6em]

\textbf{¿Qué es una cuestión federal compleja indirecta?} &
El conflicto normativo se da entre dos o más normas, todas ellas inferiores a la Constitución (por ejemplo, una norma provincial frente a una norma del Congreso, o un tratado sin jerarquía constitucional frente a una ley nacional). \\
\addlinespace[0.6em]

\textbf{¿Cuándo la decisión es contraria al derecho federal invocado?} &
Cuando el recurrente logra demostrar esa contradicción. No la hay, por ejemplo, si entre una norma provincial y una federal se resolvió que prevalece la federal, porque la Constitución ya prevé su superioridad sobre las locales. \\
\addlinespace[0.6em]

\textbf{¿Qué exige la relación directa e inmediata?} &
Que la resolución de la cuestión federal tenga relación directa con la decisión del juicio, sin tratarse de una controversia meramente conjetural o académica. \\

\midrule
\multicolumn{2}{@{}p{\dimexpr0.95\textwidth+2\tabcolsep\relax}@{}}{%
\textbf{Resumen:} Los requisitos propios son: sentencia definitiva o asimilable (con excepciones para medidas cautelares en casos graves), cuestión federal (simple, compleja directa o compleja indirecta), decisión contraria al derecho federal invocado y relación directa e inmediata entre la cuestión federal y la decisión del juicio.
} \\
\bottomrule
\end{longtable}

%==================================================================
% CAPÍTULO 4
%==================================================================
\chapter[La queja por recurso denegado]{La queja por recurso extraordinario denegado}

\section{Concepto y diferencias con el recurso extraordinario}

El recurso extraordinario no se cobra con una tasa, pero si se rechaza el recurso extraordinario, la única instancia que queda es la \textbf{queja por recurso extraordinario denegado}.

En cuanto al costo de la queja, se señaló que es elevado y que, si la queja no es aceptada, el monto depositado se destina a la biblioteca de la Corte.

\begin{table}[H]
\centering
\begin{tabularx}{\textwidth}{@{}>{\raggedright\arraybackslash}p{0.20\textwidth}>{\raggedright\arraybackslash}X>{\raggedright\arraybackslash}X@{}}
\toprule
 & \textbf{Recurso extraordinario} & \textbf{Queja} \\
\midrule
\textbf{Ante quién se presenta} &
Ante el propio tribunal que dictó el acto impugnado. &
Directamente ante la Corte Suprema. \\
\addlinespace[0.4em]
\textbf{Costo} &
No se cobra tasa. &
Requiere un depósito, salvo excepciones (por ejemplo, el trabajador). \\
\addlinespace[0.4em]
\textbf{Extensión máxima} &
40 páginas. &
10 páginas. \\
\addlinespace[0.4em]
\textbf{Cuándo procede} &
Contra la sentencia definitiva o asimilable. &
Cuando el recurso extraordinario fue rechazado. \\
\bottomrule
\end{tabularx}
\caption{Recurso extraordinario y queja}
\end{table}

Quien se encuentra en una provincia dispone de más días para presentar la queja, en razón de la distancia hasta la sede de la Corte, que se encuentra en la capital.

\section{El circuito completo del recurso extraordinario y la queja}

\begin{enumerate}
    \item Dentro de los 10 días de notificada la sentencia debe interponerse el recurso extraordinario, con una extensión máxima de 40 páginas y otros límites de renglones y de tamaño de letra.
    \item Se corre traslado del recurso extraordinario a la contraparte.
    \item Resuelve el mismo tribunal que dictó la resolución impugnada, admitiendo o no el recurso extraordinario.
    \item Si el recurso extraordinario es rechazado, puede interponerse la queja, que se presenta directamente ante la Corte.
\end{enumerate}

\section{Consideración final sobre los requisitos}

\begin{important}
Todos los requisitos del recurso extraordinario son necesarios y deben estar presentes: cualquier imperfección en alguno de ellos puede generar el rechazo de los recursos.
\end{important}

Estos requisitos permiten ir más a la sustancia del recurso y conocer sus pormenores, que no son pocos. Se recomienda repasar los tres bloques trabajados: requisitos formales, requisitos propios y requisitos comunes a todos los recursos.

\clearpage
\section{Cuadro Cornell}

\begin{longtable}{@{}>{\raggedright\arraybackslash}p{0.27\textwidth}>{\raggedright\arraybackslash}p{0.68\textwidth}@{}}
\toprule
\textbf{Preguntas / palabras clave} & \textbf{Notas / respuestas} \\
\midrule
\endfirsthead
\toprule
\textbf{Preguntas / palabras clave} & \textbf{Notas / respuestas} \\
\midrule
\endhead
\endfoot
\endlastfoot

\textbf{¿Qué es la queja por recurso extraordinario denegado?} &
Es la única instancia que queda cuando se rechaza el recurso extraordinario. Se presenta directamente ante la Corte Suprema. \\
\addlinespace[0.6em]

\textbf{¿Qué diferencia al recurso extraordinario de la queja?} &
El recurso extraordinario se presenta ante el propio tribunal que dictó el acto, no paga tasa, tiene una extensión máxima de 40 páginas y procede contra sentencia definitiva o asimilable. La queja se presenta directamente ante la Corte, requiere depósito (salvo excepciones, como el trabajador), tiene una extensión máxima de 10 páginas y procede cuando el recurso extraordinario fue rechazado. \\
\addlinespace[0.6em]

\textbf{¿Qué plazo tiene quien se encuentra en una provincia para presentar la queja?} &
Dispone de más días, en razón de la distancia hasta la sede de la Corte. \\
\addlinespace[0.6em]

\textbf{¿Cuál es el circuito completo del REF y la queja?} &
Interposición del REF dentro de los 10 días de notificada la sentencia (máximo 40 páginas, con límites de renglones y tamaño de letra); traslado a la contraparte; resolución del mismo tribunal que dictó la resolución impugnada sobre su admisión; y, si es rechazado, queja directa ante la Corte. \\
\addlinespace[0.6em]

\textbf{¿Por qué son necesarios todos los requisitos?} &
Porque cualquier imperfección en alguno de ellos puede generar el rechazo de los recursos. \\

\midrule
\multicolumn{2}{@{}p{\dimexpr0.95\textwidth+2\tabcolsep\relax}@{}}{%
\textbf{Resumen:} Si el REF es rechazado, la única vía que resta es la queja, presentada directamente ante la Corte, con depósito y extensión máxima de 10 páginas. El circuito completo va desde la interposición dentro de los 10 días ante el mismo tribunal hasta la eventual queja.
} \\
\bottomrule
\end{longtable}

%==================================================================
% CAPÍTULO 5
%==================================================================
\chapter[Control de constitucionalidad reciente]{Jurisprudencia y control de constitucionalidad reciente}

Se identificaron ejemplos recientes de cuestiones constitucionales o de declaraciones de invalidez de normas, a fin de ilustrar la aplicación de los conceptos estudiados.

\section{El Decreto de Necesidad y Urgencia 70/2023}

El DNU 70/2023 fue confirmado como ejemplo válido de cuestión constitucional reciente.

\section{La reforma laboral y la actualización de créditos judiciales}

Se mencionó, como ejemplo, la última ley laboral, respecto de la cual se habría dictado la invalidez de un artículo (identificado como el 33) relativo a las indemnizaciones, y se habría anulado la fórmula mencionada como «el 67\,\% más 3 puntos de IPC». % TODO: contenido incompleto o ambiguo en las notas originales (los datos fueron aportados con incertidumbre: "si no me equivoco")

La nueva ley establece el modo de actualización de los créditos judiciales laborales. En los últimos meses hubo una enorme litigiosidad, muy compleja: el fuero laboral, por aplicación del principio de preferente tutela del trabajador, tiende a ser más partidario del resguardo de los derechos y del patrimonio del trabajador.

Producto también del proceso inflacionario del país, habían existido mecanismos de actualización muy complejos que, en algunos casos, terminaban multiplicando varias veces el monto original de la deuda: por ejemplo, en 6 años, el 8.000 o el 10.000\,\% de la deuda. Esta disparidad se atribuye a que los distintos fueros fueron fijando, cada uno por su lado, la doctrina que consideraban más razonable.

\section{El impuesto a las ganancias sobre los haberes}

Se mencionó el impuesto a las ganancias en lo relativo a las impugnaciones con transferencia de haberes. % TODO: contenido incompleto o ambiguo en las notas originales (no se recordó si una norma fue declarada inconstitucional)
Se trata de un ámbito de mucha litigiosidad; sin embargo, lo relevante para el análisis del recurso extraordinario no es la cantidad de casos, sino verificar si se trata de la impugnación de una normativa federal.

\section{El \textit{per saltum} en la reforma laboral}

\begin{definition}{Per saltum}
Figura utilizada en el ámbito federal para pasar de la primera instancia a la Corte Suprema de forma directa, sin pasar por la cámara. Requiere acreditar una situación de enorme celeridad para resolver una cuestión de trascendencia, que permita saltear las instancias típicas del proceso.
\end{definition}

En el marco de la reforma laboral, un juez de primera instancia había declarado la inconstitucionalidad de un artículo. El gobierno solicitó el \textit{per saltum} frente a la medida cautelar que había ordenado la suspensión de varios artículos de la ley.

\begin{example}{Rechazo del per saltum solicitado por el gobierno}
La Corte rechazó la pretensión de \textit{per saltum}, al no encontrar acreditada la extrema celeridad invocada, y consideró razonable que la causa siguiera su curso ordinario ante la cámara. La cámara terminó levantando la suspensión de los artículos de la ley que había dispuesto la medida cautelar de primera instancia.
\end{example}

\section{Límites de la Corte a los DNU: la ley de lealtad comercial}

Se describió un caso vinculado a decretos de necesidad y urgencia dictados en 2019, que modificaron por decreto la ley de lealtad comercial, y que constituye un ejemplo relevante y actual de control de constitucionalidad.

\begin{definition}{Ley de lealtad comercial}
Ley que prevé determinadas conductas consideradas violatorias del juego limpio o de la lealtad comercial que debe existir entre empresas que compiten en un mismo mercado, sancionando esas conductas como infracciones.
\end{definition}

La Corte declaró que, en materia de infracciones, no pueden fijarse por decreto determinadas cuestiones que son claves para la determinación de una multa: debe ser una ley del Congreso.

\clearpage
\section{Cuadro Cornell}

\begin{longtable}{@{}>{\raggedright\arraybackslash}p{0.27\textwidth}>{\raggedright\arraybackslash}p{0.68\textwidth}@{}}
\toprule
\textbf{Preguntas / palabras clave} & \textbf{Notas / respuestas} \\
\midrule
\endfirsthead
\toprule
\textbf{Preguntas / palabras clave} & \textbf{Notas / respuestas} \\
\midrule
\endhead
\endfoot
\endlastfoot

\textbf{¿Qué ejemplos recientes de control de constitucionalidad se identificaron?} &
El DNU 70/2023; la reforma laboral y la actualización de créditos judiciales; el impuesto a las ganancias sobre haberes (mencionado sin precisión); el \textit{per saltum} solicitado en la reforma laboral; y los DNU de 2019 que modificaron la ley de lealtad comercial. \\
\addlinespace[0.6em]

\textbf{¿Por qué generó litigiosidad la actualización de créditos judiciales laborales?} &
Porque, por el proceso inflacionario, existieron mecanismos de actualización muy complejos que llegaron a multiplicar varias veces el monto de la deuda (8.000 o 10.000\,\% en 6 años), y porque cada fuero fue fijando por su lado la doctrina que consideraba más razonable. \\
\addlinespace[0.6em]

\textbf{¿Qué es el \textit{per saltum}?} &
Figura que permite pasar de la primera instancia a la Corte Suprema de forma directa, sin pasar por la cámara, cuando se acredita una situación de enorme celeridad para resolver una cuestión de trascendencia. \\
\addlinespace[0.6em]

\textbf{¿Qué ocurrió con el \textit{per saltum} solicitado en la reforma laboral?} &
El gobierno lo solicitó frente a la cautelar que suspendía varios artículos de la ley. La Corte lo rechazó por no estar acreditada la extrema celeridad y consideró razonable que la causa siguiera ante la cámara, que terminó levantando la suspensión. \\
\addlinespace[0.6em]

\textbf{¿Qué límite puso la Corte a los DNU en materia de lealtad comercial?} &
Declaró que en materia de infracciones no pueden fijarse por decreto cuestiones clave para determinar una multa: debe ser una ley del Congreso. \\

\midrule
\multicolumn{2}{@{}p{\dimexpr0.95\textwidth+2\tabcolsep\relax}@{}}{%
\textbf{Resumen:} Los ejemplos recientes (DNU 70/2023, reforma laboral, \textit{per saltum} rechazado y DNU sobre lealtad comercial) muestran cómo se aplica el control de constitucionalidad. La Corte rechazó el \textit{per saltum} por falta de extrema celeridad y exigió una ley del Congreso para fijar por decreto cuestiones clave en materia de multas.
} \\
\bottomrule
\end{longtable}

%==================================================================
% CAPÍTULO 6
%==================================================================
\chapter[Síntesis final]{Síntesis final}

El recurso extraordinario federal es la vía procesal que permite llevar una causa hasta la Corte Suprema cuando está comprometida la supremacía de la Constitución. Su procedencia depende del cumplimiento estricto de tres tipos de requisitos, que se resumen en el siguiente cuadro.

\begin{table}[H]
\centering
\begin{tabularx}{\textwidth}{@{}>{\raggedright\arraybackslash}p{0.22\textwidth}>{\raggedright\arraybackslash}X@{}}
\toprule
\textbf{Tipo de requisito} & \textbf{Contenido} \\
\midrule
Formales &
Regulados en gran medida por la Acordada 4/2007: extensión, plazo de 10 días, tribunal de interposición y carátula. Incluyen el planteo oportuno del caso federal en cada instancia. \\
\addlinespace[0.4em]
Comunes a todo recurso &
Gravamen actual y existencia de un proceso judicial. \\
\addlinespace[0.4em]
Propios &
Núcleo sustantivo de la figura: sentencia definitiva o asimilable, cuestión federal (simple, o compleja directa e indirecta), decisión contraria al derecho federal invocado y relación directa e inmediata entre esa cuestión y el fallo. \\
\bottomrule
\end{tabularx}
\caption{Requisitos del recurso extraordinario federal}
\end{table}

Si el recurso es rechazado, la única vía que resta es la queja, presentada directamente ante la Corte.

Los ejemplos jurisprudenciales trabajados (el DNU 70/2023, la reforma laboral y el rechazo del \textit{per saltum}, el precedente sobre la ley de lealtad comercial, el caso del Grupo Clarín frente a la Ley de Medios y la disputa por el financiamiento universitario) muestran que estos requisitos no son un formalismo vacío, sino el filtro que define, en la práctica, qué controversias llegan a ser resueltas por el máximo tribunal del país.

\end{document}
